
\documentclass[11pt,a4paper]{article}

\usepackage{kotex}
\usepackage{amsmath,amssymb,amsfonts}
\usepackage{bm}
\usepackage{geometry}
\usepackage{graphicx}
\usepackage{booktabs}
\usepackage{algorithm}
\usepackage{algpseudocode}
\usepackage{hyperref}

\geometry{margin=25mm}

\title{
Support Tensor Theory for Additive Manufacturing Orientation Optimization\\
\large 고유벡터 기반 3D 프린팅 지지구조량 예측 이론
}

\author{
Research Draft
}

\date{\today}

\begin{document}

\maketitle

\begin{abstract}
In additive manufacturing, the amount of support structure strongly depends on the build orientation of a three-dimensional object.
Conventional orientation optimization often requires repeated support-volume simulations over many sampled directions.
This paper proposes a compact tensor-based approximation, called the \emph{Support Tensor}, that estimates favorable build directions directly from the surface normal distribution of a triangular mesh.
For a mesh composed of triangular faces with area $A_i$ and unit outward normal $\bm m_i$, the Support Tensor is defined as
\[
M = \sum_{i=1}^{N} A_i \bm m_i \bm m_i^T.
\]
The directional support tendency is approximated by the quadratic form
\[
\widetilde S(\bm n)=\bm n^T M \bm n,
\]
where $\bm n$ is the unit build direction.
By the Rayleigh quotient theorem, the eigenvector associated with the smallest eigenvalue of $M$ gives a candidate direction that minimizes the approximate support tendency.
The framework is further extended to cavity-aware or overhang-aware weighting by introducing face weights $w_i$.
The proposed method provides a fast, interpretable, and mathematically tractable preprocessing tool for build-orientation selection in 3D printing.
\end{abstract}

\section{Introduction}

In material extrusion, vat photopolymerization, powder bed fusion, and other additive manufacturing processes, the build orientation of a part significantly affects support volume, surface quality, build time, material consumption, and post-processing cost.
A common engineering problem is therefore to find a build direction that minimizes the amount of support structure while satisfying other manufacturing constraints.

Let $\bm n \in \mathbb R^3$ be a unit vector representing the build direction.
The ideal support volume may be described as a direction-dependent functional
\[
S(\bm n),
\qquad
\|\bm n\|=1.
\]
The orientation optimization problem is
\begin{equation}
\min_{\|\bm n\|=1} S(\bm n).
\label{eq:main_problem}
\end{equation}

Direct evaluation of $S(\bm n)$ is often computationally expensive because it may require slicing, ray casting, visibility checks, overhang detection, and support generation.
This motivates a fast approximation that can identify promising candidate directions before detailed simulation.

The main idea of this paper is to construct a $3\times 3$ symmetric matrix from the surface normal distribution of the object.
The eigenvectors of this matrix then provide physically interpretable candidate build directions.
This approach is analogous to principal component analysis, but instead of using point coordinates, it uses surface normals weighted by triangle areas.

\section{Triangular Mesh Model}

Assume that a solid object is represented by a closed triangular surface mesh
\[
\mathcal M = \{T_i\}_{i=1}^{N}.
\]
For each triangle $T_i$, define

\begin{itemize}
\item $A_i > 0$: area of the triangle,
\item $\bm m_i \in \mathbb R^3$: unit outward normal vector,
\item $\bm c_i \in \mathbb R^3$: centroid of the triangle.
\end{itemize}

The build direction is denoted by a unit vector
\[
\bm n \in \mathbb S^2
=
\{\bm n\in\mathbb R^3:\|\bm n\|=1\}.
\]

\section{Support Functional Based on Overhang}

A face tends to require support when its outward normal points downward relative to the build direction.
A simple overhang indicator can be written as
\[
\bm m_i\cdot \bm n < 0.
\]
The negative part of a scalar $x$ is defined as
\[
x_-=\max(0,-x).
\]
Using this notation, a first-order support tendency functional can be written as
\begin{equation}
S_1(\bm n)
=
\sum_{i=1}^{N} A_i (\bm m_i\cdot \bm n)_-.
\label{eq:first_order_support}
\end{equation}

Equivalently, using the Heaviside function $H$,
\begin{equation}
S_1(\bm n)
=
\sum_{i=1}^{N}
A_i
H(-\bm m_i\cdot \bm n)
(-\bm m_i\cdot \bm n).
\label{eq:heaviside_support}
\end{equation}

This expression captures the intuitive fact that large downward-facing surface areas tend to require more support.
However, because of the Heaviside function and the negative-part operator, $S_1(\bm n)$ is not a simple quadratic function of $\bm n$.

\section{Quadratic Approximation}

To obtain an eigenvalue-based formulation, we approximate the directional dependence by a quadratic form.
Consider
\begin{equation}
\widetilde S(\bm n)
=
\sum_{i=1}^{N}
A_i
(\bm m_i\cdot \bm n)^2.
\label{eq:quadratic_support}
\end{equation}

Since
\[
(\bm m_i\cdot \bm n)^2
=
\bm n^T
(\bm m_i\bm m_i^T)
\bm n,
\]
we have
\begin{align}
\widetilde S(\bm n)
&=
\sum_{i=1}^{N}
A_i
\bm n^T
(\bm m_i\bm m_i^T)
\bm n\\
&=
\bm n^T
\left(
\sum_{i=1}^{N}
A_i
\bm m_i\bm m_i^T
\right)
\bm n.
\end{align}

This motivates the following definition.

\section{Definition of the Support Tensor}

\textbf{Definition 1.}
For a triangular mesh $\mathcal M$, the area-weighted Support Tensor is defined as
\begin{equation}
M
=
\sum_{i=1}^{N}
A_i
\bm m_i\bm m_i^T.
\label{eq:support_tensor}
\end{equation}

The corresponding quadratic support tendency is
\begin{equation}
\widetilde S(\bm n)
=
\bm n^T M \bm n.
\label{eq:support_quadratic_form}
\end{equation}

The matrix $M$ is symmetric and positive semidefinite because each term
$\bm m_i\bm m_i^T$ is symmetric and positive semidefinite.
Therefore, all eigenvalues of $M$ are real and nonnegative.

\section{Eigenvalue Formulation}

The optimization problem associated with the quadratic approximation is
\begin{equation}
\min_{\|\bm n\|=1}
\bm n^T M \bm n.
\label{eq:rayleigh_problem}
\end{equation}

Let the eigenvalues of $M$ be
\[
\lambda_1 \leq \lambda_2 \leq \lambda_3,
\]
with corresponding orthonormal eigenvectors
\[
\bm v_1,\bm v_2,\bm v_3.
\]

By the Rayleigh quotient theorem,
\begin{equation}
\min_{\|\bm n\|=1}
\bm n^T M \bm n
=
\lambda_1,
\end{equation}
and the minimum is attained at
\begin{equation}
\bm n = \bm v_1.
\end{equation}

Thus, the eigenvector corresponding to the smallest eigenvalue provides a candidate build direction that minimizes the approximate support tendency.

\section{Relation to Principal Component Analysis}

Principal component analysis constructs a covariance matrix from point coordinates:
\begin{equation}
C
=
\sum_{j=1}^{P}
(\bm x_j-\bar{\bm x})
(\bm x_j-\bar{\bm x})^T.
\end{equation}

The eigenvectors of $C$ represent principal geometric axes of the point cloud.
In contrast, the Support Tensor uses surface normals:
\begin{equation}
M
=
\sum_{i=1}^{N}
A_i
\bm m_i\bm m_i^T.
\end{equation}

Therefore, PCA describes the distribution of mass or geometry in space, while the Support Tensor describes the distribution of surface orientations.
For support-structure prediction, the surface normal distribution is often more directly relevant than the point distribution.

\section{Weighted Support Tensor}

The basic Support Tensor can be generalized by assigning a nonnegative weight $w_i$ to each triangle:
\begin{equation}
M_w
=
\sum_{i=1}^{N}
w_i A_i
\bm m_i\bm m_i^T.
\label{eq:weighted_support_tensor}
\end{equation}

The corresponding directional functional is
\begin{equation}
\widetilde S_w(\bm n)
=
\bm n^T M_w \bm n.
\end{equation}

The weight $w_i$ can encode additional manufacturing or geometric information, such as

\begin{itemize}
\item overhang sensitivity,
\item local surface roughness requirement,
\item distance from the build plate,
\item internal cavity influence,
\item accessibility for support removal,
\item material-dependent support penalty.
\end{itemize}

\subsection{Cavity-Aware Weighting}

For a part with internal concave spaces or cavities, support difficulty may depend not only on surface orientation but also on whether a region is internally trapped or difficult to access.
A simple cavity-aware model can be written as
\begin{equation}
w_i
=
1+\alpha q_i,
\end{equation}
where

\begin{itemize}
\item $\alpha\geq 0$ is a tuning parameter,
\item $q_i\in[0,1]$ is a cavity influence score for triangle $i$.
\end{itemize}

Then
\begin{equation}
M_{\text{cav}}
=
\sum_{i=1}^{N}
(1+\alpha q_i)
A_i
\bm m_i\bm m_i^T.
\end{equation}

If $q_i$ is large for faces near internal cavities or concave pockets, the resulting eigenvectors become more sensitive to cavity-related support penalties.

\section{Comparison with Buoyancy Analogy}

The proposed idea is related to a buoyancy analogy.
Consider an ice-like body containing internal air cavities.
The orientation of the body in water depends on the relationship between the center of mass and the center of buoyancy.
If the object has internal voids, the distribution of mass and displaced volume may differ, producing orientation-dependent torques.

Similarly, in 3D printing, concave spaces and downward-facing surfaces generate orientation-dependent support requirements.
Although exact buoyancy or support behavior is generally nonlinear, a tensor approximation can capture first-order directional tendencies.
The Support Tensor is therefore not an exact physical law, but a compact directional descriptor useful for fast orientation screening.

\section{Algorithm}

\begin{algorithm}
\caption{Support Tensor-Based Build Direction Estimation}
\begin{algorithmic}[1]
\Require Triangular mesh $\mathcal M=\{T_i\}_{i=1}^{N}$
\Ensure Candidate build direction $\bm n^\ast$
\State Initialize $M\leftarrow 0_{3\times 3}$
\For{$i=1$ to $N$}
    \State Compute triangle area $A_i$
    \State Compute unit outward normal $\bm m_i$
    \State Choose weight $w_i$; use $w_i=1$ for the basic model
    \State Update $M\leftarrow M+w_i A_i\bm m_i\bm m_i^T$
\EndFor
\State Compute eigenvalues and eigenvectors of $M$
\State Let $\bm v_1$ be the eigenvector corresponding to the smallest eigenvalue
\State Set $\bm n^\ast\leftarrow \bm v_1$
\State Optionally evaluate both $\bm v_1$ and $-\bm v_1$ using the original support functional
\State \Return $\bm n^\ast$
\end{algorithmic}
\end{algorithm}

\section{Post-Selection Using the Original Functional}

Because the quadratic form satisfies
\[
\widetilde S(\bm n)=\widetilde S(-\bm n),
\]
it cannot distinguish upward and downward directions.
However, the actual support functional
\[
S_1(\bm n)
=
\sum_{i=1}^{N}
A_i
(\bm m_i\cdot \bm n)_-
\]
is generally not identical for $\bm n$ and $-\bm n$.
Therefore, after computing the eigenvector $\bm v_1$, one should evaluate
\[
S_1(\bm v_1)
\quad\text{and}\quad
S_1(-\bm v_1)
\]
and choose the better of the two.

A practical decision rule is
\begin{equation}
\bm n^\ast
=
\arg\min_{\bm n\in\{\bm v_1,-\bm v_1\}}
S_1(\bm n).
\end{equation}

\section{Experimental Design}

To validate the proposed method, one may compare the following approaches:

\begin{enumerate}
\item exhaustive directional sampling on the unit sphere,
\item PCA-based build direction estimation,
\item basic Support Tensor eigenvector method,
\item weighted or cavity-aware Support Tensor method,
\item commercial slicer-based support volume estimation.
\end{enumerate}

For each mesh, compute

\begin{itemize}
\item predicted build direction,
\item estimated support area or volume,
\item computation time,
\item angular error relative to the best sampled direction,
\item material saving ratio.
\end{itemize}

A possible performance metric is
\begin{equation}
E_{\text{angle}}
=
\cos^{-1}
\left(
|\bm n_{\text{pred}}\cdot \bm n_{\text{best}}|
\right).
\end{equation}

Another metric is the relative support penalty:
\begin{equation}
E_{\text{support}}
=
\frac{
S(\bm n_{\text{pred}})-S(\bm n_{\text{best}})
}{
S(\bm n_{\text{best}})
}.
\end{equation}

\section{Discussion}

The Support Tensor method has several advantages.

First, it is computationally inexpensive.
The main computation is the construction and eigenvalue decomposition of a $3\times 3$ symmetric matrix.

Second, it is interpretable.
Each triangle contributes according to its area and normal direction, so the tensor directly reflects the surface orientation distribution.

Third, it can be extended.
Weights can encode overhang angle thresholds, cavity difficulty, surface quality constraints, or material-dependent penalties.

However, the method also has limitations.
The quadratic approximation does not fully represent actual support volume, which depends on visibility, connectivity, slicing, self-shadowing, support removal, and process-specific rules.
Therefore, the method should be regarded as a fast candidate-generation technique rather than a complete replacement for slicer-based simulation.

\section{Conclusion}

This paper proposed a tensor-based method for estimating favorable build directions in additive manufacturing.
By defining the Support Tensor
\[
M
=
\sum_{i=1}^{N}
A_i
\bm m_i\bm m_i^T,
\]
the approximate support tendency becomes a quadratic form
\[
\widetilde S(\bm n)=\bm n^T M\bm n.
\]
The eigenvector associated with the smallest eigenvalue of $M$ gives a mathematically justified candidate direction for minimizing support tendency.
The method is simple, fast, and extensible to cavity-aware and process-aware weighting schemes.
It may serve as a useful preprocessing step for build-orientation optimization and support-structure reduction in 3D printing.

\section*{Future Work}

Future research may include

\begin{itemize}
\item deriving better approximations from the nonlinear support functional,
\item incorporating overhang threshold angles,
\item modeling internal cavities using visibility or ray-casting scores,
\item comparing the method with commercial slicer outputs,
\item validating the approach on thermal manikin and functional garment-related 3D printed models,
\item extending the tensor model to multi-objective optimization including strength, surface quality, and build time.
\end{itemize}

\bibliographystyle{plain}
\begin{thebibliography}{9}

\bibitem{pca}
Jolliffe, I. T.
\textit{Principal Component Analysis}.
Springer, 2002.

\bibitem{am}
Gibson, I., Rosen, D., and Stucker, B.
\textit{Additive Manufacturing Technologies}.
Springer, 2015.

\bibitem{tensor}
Mardia, K. V. and Jupp, P. E.
\textit{Directional Statistics}.
Wiley, 2000.

\end{thebibliography}

\end{document}
